\SourcePage{307}{312}
\section{Invariant amplitudes}

Helicity amplitudes use a particular reference frame, namely the center-of-momentum frame. When scattering amplitudes are calculated by invariant perturbation theory (and also when their general analytic properties are studied), however, it is convenient to write them in a manifestly invariant form.

If the particles participating in the reaction have no spin, the scattering amplitude depends only on invariant products of their four-momenta. For a reaction of the form
\begin{equation}
a+b\to c+d
\tag{70.1}
\end{equation}
any two of the quantities $s$, $t$, and $u$ defined in \S~66 may be chosen as these invariants. The scattering amplitude then reduces to a single function, $M_{fi}=f(s,t)$.

If the particles have spin, there are, in addition to the kinematic invariants $s$, $t$, and $u$, invariants that can be formed from the particle wave amplitudes (bispinors, four-tensors, and so on). The scattering amplitudes must then have the form
\begin{equation}
M_{fi}=\sum_n f_n(s,t)F_n,
\tag{70.2}
\end{equation}
where the $F_n$ are invariants depending linearly on the wave amplitudes of all participating particles (and also on their four-momenta). The coefficients $f_n(s,t)$ are called \emph{invariant amplitudes}.

If the wave amplitudes are chosen to describe particles of definite helicity, the invariants acquire definite values $F_n=F_n(\lambda_i,\lambda_f)$. The helicity scattering amplitudes are then homogeneous linear combinations of the invariant amplitudes $f_n$. It follows that the number of independent functions $f_n(s,t)$ equals the number of independent helicity amplitudes. Since the latter number is readily found (as explained in \S~69), the construction of the invariants $F_n$ is thereby simplified: their required number is known in advance.

Let us consider some examples. In every example we assume that the interaction is $T$- and $P$-invariant; the latter
\SourcePage{308}{313}
property means that the invariants $F_n$ must be true scalars rather than pseudoscalars.

\subsection*{Scattering of a spin-0 particle by a spin-$1/2$ particle.}

To count the invariants---or, equivalently, the independent helicity amplitudes---note that the total number of elements of the matrix $S^J$ (that is, the number of different sets of values $\lambda_1$, $\lambda_2$, $\lambda_1'$, $\lambda_2'$) is 4 in this case ($\lambda_1=\lambda_1'=0$, $\lambda_2$, $\lambda_2'=\pm1/2$). Taking $P$ invariance into account reduces the number of independent elements to two, after which $T$ invariance produces no further reduction.

Two independent invariants may be chosen as
\begin{equation}
F_1=\bar u'u,\qquad F_2=\bar u'(\gamma K)u.
\tag{70.3}
\end{equation}
Here $u=u(p)$ and $u'=u(p')$ are the bispinor amplitudes of the initial and final fermions; $K=k+k'$, where $k$ and $k'$ are the four-momenta of the initial and final bosons\footnote{At first sight, one might also form an invariant of the type $\bar u'\sigma_{\mu\nu}k^\mu k'^{\nu}u$ (the matrices $\sigma_{\mu\nu}$ are defined in (28.2)). It is readily shown, however, that this invariant reduces to those in (70.3) when the conservation law $k'=p+k-p'$ and the equations
\[
(\gamma p)u=mu,\qquad \bar u'(\gamma p')=m\bar u',
\]
satisfied by the bispinor amplitudes are used.}.

The $T$ invariance of the quantities (70.3) becomes evident on noting that the products $\bar u'u$ and $\bar u'\gamma^\mu u$ transform under time reversal by the same law (28.6) as the operators $\widehat{\bar\psi\psi}$ and $\widehat{\bar\psi\gamma^\mu\psi}$ of which they are matrix elements: the product $\bar u'u$ is invariant by itself, whereas the four-vector $\bar u'\gamma u$ transforms as
\[
\bar u'\gamma^0u\to\bar u'\gamma^0u,
\qquad
\bar u'\boldsymbol\gamma u\to-\bar u'\boldsymbol\gamma u.
\]
The four-momenta transform in the same way, $(K^0,\mathbf K)\to(k^0,-\mathbf K)$, and consequently the scalar product $F_2=K_\mu(\bar u'\gamma^\mu u)$ is invariant.

\subsection*{Elastic scattering of two identical spin-$1/2$ particles.}

To count the independent helicity amplitudes, it is convenient to begin with the linear combinations of helicity states
\[
\psi_{1g}=\psi_{++}+\psi_{--},\qquad
\psi_{2g}=\psi_{++}-\psi_{--},
\]
\[
\psi_{3g}=\psi_{+-}+\psi_{-+},\qquad
\psi_u=\psi_{+-}-\psi_{-+},
\]
where the subscripts $+$ and $-$ indicate the helicities ($\pm1/2$) of the two particles. The states $1g$, $2g$, and $3g$ are even, while the state $u$ is odd under particle interchange. The transitions $g\leftrightarrow u$ are therefore forbidden, leaving $16-6=10$ matrix elements after interchange symmetry is taken into account. Under the inversion
\SourcePage{309}{314}
$P$, the functions $\psi_{1g}$, $\psi_{3g}$ have parity opposite to that of $\psi_{2g}$; forbidding transitions between them reduces the number of independent amplitudes to six. Finally, $T$ invariance makes the amplitudes of the transitions $1g\to3g$ and $3g\to1g$ equal, so that only five independent amplitudes remain. Five independent invariants may be chosen as
\begin{equation}
\begin{aligned}
F_1&=(\bar u_1'u_1)(\bar u_2'u_2),
&F_2&=(\bar u_1'\gamma^5u_1)(\bar u_2'\gamma^5u_2),\\
F_3&=(\bar u_1'\gamma^\mu u_1)(\bar u_2'\gamma_\mu u_2),
&F_4&=(\bar u_1'\gamma^\mu\gamma^5u_1)(\bar u_2'\gamma_\mu\gamma^5u_2),\\
F_5&=(\bar u_1'\sigma^{\mu\nu}u_1)(\bar u_2'\sigma_{\mu\nu}u_2),
\end{aligned}
\tag{70.4}
\end{equation}
where $u_1$, $u_2$ are the bispinor amplitudes of the initial particles, and $u_1'$, $u_2'$ those of the final particles. Interchanging the initial (or final) particles produces no new invariants: the resulting invariants can be expressed through the original ones (see the problem to \S~28). But expression (70.2) with the $F_n$ from (70.4) does not explicitly incorporate the requirement that interchanging two identical fermions must change the sign of the scattering amplitude. An expression satisfying this requirement can be written as
\begin{equation}
M_{fi}=\bigl[(\bar u_1'u_1)(\bar u_2'u_2)f_1(t,u)
-(\bar u_2'u_1)(\bar u_1'u_2)f_1(u,t)\bigr]+\ldots
\tag{70.5}
\end{equation}
When $p_1'$ and $p_2'$ (or $p_1$ and $p_2$) are interchanged, the kinematic invariants transform as $s\to s$, $t\to u$, $u\to t$, so the stated requirement is fulfilled automatically.

\subsection*{Elastic scattering of a photon by particles of spin 0 and $1/2$.}

It is useful to express the amplitudes of these processes in terms of unit spacelike four-vectors $e^{(1)}$, $e^{(2)}$ satisfying
\begin{equation}
\begin{aligned}
e^{(1)2}=e^{(2)2}&=-1,&e^{(1)}e^{(2)}&=0,\\
e^{(1)}k=e^{(2)}k&=0,&e^{(1)}k'=e^{(2)}k'&=0
\end{aligned}
\tag{70.6}
\end{equation}
(for each of the two photons these four-vectors may serve as the four-unit vectors used in an invariant description of its polarization properties; see \S~8).

Let $k$ and $k'$ be the initial and final photon four-momenta, and let $p$ and $p'$ be the corresponding four-momenta of the scattering particle. Consider the four-vectors
\begin{equation}
P^\lambda=p^\lambda+p'^\lambda
-K^\lambda\frac{pK+p'K}{K^2},
\qquad
N^\lambda=e^{\lambda\mu\nu\rho}P_\mu q_\nu K_\rho,
\tag{70.7}
\end{equation}
where
\[
K=k+k',\qquad q=p-p'=k'-k.
\]
They are manifestly mutually orthogonal. They are also orthogonal to the four-vectors $K$, $q$, and hence to $k$, $k'$. Since they are orthogonal to the timelike four-vector $K$ ($K^2=2kk'>0$), they are themselves spacelike (indeed, in a frame where $\mathbf K=0$, the relation $KP=0$ gives $P_0=0$, and therefore
\SourcePage{310}{315}
$P^2<0$). Normalizing $P$ and $N$, that is, forming
\begin{equation}
e^{(1)\lambda}=\frac{N^\lambda}{\sqrt{-N^2}},
\qquad
e^{(2)\lambda}=\frac{P^\lambda}{\sqrt{-P^2}},
\tag{70.8}
\end{equation}
gives a pair of four-vectors with all the required properties. Notice that $e^{(2)}$ is a true vector, whereas $e^{(1)}$ is a pseudovector. Write the photon scattering amplitude as
\begin{equation}
M_{fi}=F^{\lambda\mu}e_\lambda^{\prime *}e_\mu,
\tag{70.9}
\end{equation}
with the polarization four-vectors $e$ and $e'$ of the initial and final photons displayed explicitly.

Photon helicity takes only two values ($\pm1$). Hence, for the scattering of a photon by a spin-0 particle, the number of independent helicity amplitudes is the same as for the mutual scattering of particles of spin 0 and $1/2$, namely 2. The tensor $F^{\lambda\mu}$ in (70.9) must be constructed only from the particle four-momenta. It can be written as
\begin{equation}
F^{\lambda\mu}
=f_1e^{(1)\lambda}e^{(1)\mu}
+f_2e^{(2)\lambda}e^{(2)\mu},
\tag{70.10}
\end{equation}
where $f_1$, $f_2$ are invariant amplitudes. Note that $F^{\lambda\mu}$ cannot contain a term involving the product $e^{(1)\lambda}e^{(2)\mu}$, because this product is a pseudotensor and would give a pseudoscalar upon substitution into (70.9).

Finally, consider the scattering of a photon by a spin-$1/2$ particle. To count the independent helicity amplitudes, note that the total number of elements of the matrix $S^J$ is 16 in this case (the helicity of each of the two initial and two final particles takes two values). The requirement of $P$ invariance reduces this number to 8, and the requirement of $T$ invariance then reduces it to 6.

In this case, write the tensor $F_{\lambda\mu}$ as
\begin{equation}
\begin{aligned}
F_{\lambda\mu}={}&G_0\bigl(e_\lambda^{(1)}e_\mu^{(1)}
+e_\lambda^{(2)}e_\mu^{(2)}\bigr)
+G_1\bigl(e_\lambda^{(1)}e_\mu^{(2)}
+e_\lambda^{(2)}e_\mu^{(1)}\bigr)+\\
&+G_2\bigl(e_\lambda^{(1)}e_\mu^{(2)}
-e_\lambda^{(2)}e_\mu^{(1)}\bigr)
+G_3\bigl(e_\lambda^{(1)}e_\mu^{(1)}
-e_\lambda^{(2)}e_\mu^{(2)}\bigr),
\end{aligned}
\tag{70.11}
\end{equation}
where $G_0$, $G_3$ are true scalars, while $G_1$, $G_2$ are pseudoscalars. Both kinds are bilinear in the fermion bispinor amplitudes $\bar u(p')$ and $u(p)$, that is, they have the form
\begin{equation}
G_n=\bar u(p')Q_nu(p).
\tag{70.12}
\end{equation}
The general form of the matrices (in bispinor indices) $Q_n$ is
\begin{equation}
\begin{aligned}
Q_0&=f_1+f_2(\gamma K),
&Q_1&=\gamma^5\bigl[f_3+f_4(\gamma K)\bigr],\\
Q_2&=\gamma^5\bigl[f_5+f_6(\gamma K)\bigr],
&Q_3&=f_7+f_8(\gamma K),
\end{aligned}
\tag{70.13}
\end{equation}
where $K=k+k'$. The coefficients $f_1$, \ldots, $f_8$ are invariant amplitudes, whose number here is 8 (instead of the required 6) because the requirement of $T$ invariance has not yet been imposed.
\SourcePage{311}{316}

Time reversal interchanges the initial and final particle four-momenta and also changes the signs of their spatial components:
\begin{equation}
(k_0,\mathbf k)\leftrightarrow(k_0',-\mathbf k'),
\qquad
(p_0,\mathbf p)\leftrightarrow(p_0',-\mathbf p');
\tag{70.14}
\end{equation}
the photon polarization four-vectors transform according to
\begin{equation}
(e_0,\mathbf e)\leftrightarrow(e_0^{\prime *},-\mathbf e^{\prime *})
\tag{70.15}
\end{equation}
(cf. (8.11a)), so that
\[
(e_0^{\prime *}e_0,e_i^{\prime *}e_0,
e_i^{\prime *}e_k)\to
(e_0^{\prime *}e_0,-e_0^{\prime *}e_i,
e_k^{\prime *}e_i).
\]
By this last transformation, the condition that the scattering amplitude (70.9) be invariant is equivalent to the requirement
\[
(F_{00},F_{i0},F_{ik})\to(F_{00},-F_{0i},F_{ki}).
\]
On the other hand, the substitutions (70.14) give
\[
(K_0,\mathbf K)\to(K_0,-\mathbf K),
\qquad(q_0,\mathbf q)\to(-q_0,\mathbf q),
\]
\[
(P_0,\mathbf P)\to(P_0,-\mathbf P),
\qquad(N_0,\mathbf N)\to(N_0,-\mathbf N),
\]
and hence
\begin{equation}
(e_0^{(1,2)},\mathbf e^{(1,2)})
\to(e_0^{(1,2)},-\mathbf e^{(1,2)}).
\tag{70.16}
\end{equation}
It follows from (70.11) that
\[
G_{0,1,3}\to G_{0,1,3},\qquad G_2\to-G_2.
\]
But under time reversal
\[
\bar u'\gamma^5u\to-\bar u'\gamma^5u,
\qquad
\bar u'\gamma^5(\gamma K)u\to
\bar u'\gamma^5(\gamma K)u,
\]
as follows from the transformation laws for the pseudoscalar and pseudovector bilinear forms in (28.6). It is therefore evident from (70.12) and (70.13) that $T$ invariance of the scattering amplitude requires
\begin{equation}
f_3=f_6=0.
\tag{70.17}
\end{equation}
